\documentclass[10pt,UTF8]{article}

\usepackage{geometry}

\usepackage{ctex}

\geometry{a4paper,scale=0.9}
\ctexset{today=old}

\usepackage[bookmarksnumbered=true]{hyperref}  % 生成大纲


\usepackage{times}  % 设置正文字体为 Adobe Times Roman

% \usepackage{makecell}
% \usepackage{multirow} % 表格多行合并

% \usepackage{graphicx} %插入图片的宏包
% \usepackage{subfigure}
% \usepackage{float} %设置图片浮动位置的宏包

\usepackage{amsmath}
\usepackage{amssymb}
\usepackage{mathtools}

\usepackage{physics}

% \usepackage{siunitx}
% % 关闭警告
% \ExplSyntaxOn
% \msg_redirect_name:nnn { siunitx } { physics-pkg } { none }
% \ExplSyntaxOff

%%%%%%%%%%%%%%%%%%%%%%%% tcolorbox %%%%%%%%%%%%%%%%%%%%%%%%%%
\usepackage[most]{tcolorbox}

\newenvironment{definition box}
{\begin{tcolorbox}[colback=blue!5!white,colframe=blue!75!black,parbox=false,before upper=\par,before lower=\par]}
{\end{tcolorbox}}

\newenvironment{theorem box}
{\begin{tcolorbox}[colback=yellow!5!white,colframe=yellow!75!black,parbox=false,before upper=\par,before lower=\par]}
{\end{tcolorbox}}
%%%%%%%%%%%%%%%%%%%%%%%%%%%%%%%%%%%%%%%%%%%%%%%%%%%%%%%%%%%%
            
%%%%%%%%%%%%%%%%%%%%% amsthm %%%%%%%%%%%%%%%%
\usepackage{amsthm}

\theoremstyle{definition}
\newtheorem{definition inner}{Definition}[section]
\newenvironment{definition}[1][]
{\begin{definition box}\begin{definition inner}[#1]\hfill\par}
{\end{definition inner}\end{definition box}}

\theoremstyle{plain}
\newtheorem{theorem inner}{Theorem}[section]
\newenvironment{theorem}[1][]
{\begin{theorem box}\begin{theorem inner}[#1]\hfill\par}
{\end{theorem inner}\end{theorem box}}

\theoremstyle{definition}
\newtheorem*{proof inner}{\textit{Proof}}
\renewenvironment{proof}[1][]
{\begin{proof inner}[#1]\hfill\par}
{\qed\end{proof inner}}

\theoremstyle{remark}
\newtheorem*{remark}{\textit{Remark}}
%%%%%%%%%%%%%%%%%%%%%%%%%%%%%%%%%%%%%%%%%%%%%%


\newcommand{\mycases}[2][0.8]{
    \left\{\vphantom{\scalebox{#1}{
        $\begin{matrix}#2\end{matrix}$
    }}\right.\!\!\!\!
    \begin{array}{l@{\quad}l}#2\end{array}
}

\newcommand{\ju}{\mathrm{j}{\mkern 1 mu}}  % 虚数单位



\begin{document}


\section{Impulse Representation}

\paragraph*{(Sifting Property)}

The sifting property is the ability of the Dirac delta function to ``pick out'' the value of a function at a specific point during integration.

\noindent
\begin{minipage}[t]{.5\textwidth}
    \subsection{DT}
    \begin{equation*}
        x[n] = \sum_{k=-\infty}^{\infty} x[k] \delta[n-k]
        = x[n] \ast \delta[n]
    \end{equation*}
\end{minipage}\hfill
\begin{minipage}[t]{.5\textwidth}
    \subsection{CT}
    \begin{equation*}
        x(t) = \int_{-\infty}^{\infty} x(\tau) \delta(t-\tau) \dd{\tau}
        = x(t) \ast \delta(t)
    \end{equation*}
\end{minipage}

\section{Impulse Response}

Let $S$ be an \emph{LTI system} and $x[n] \xlongrightarrow{S} y[n]$.

Suppose $\delta[n] \xlongrightarrow{S} h[n]$ and $\delta[n-k]\xlongrightarrow{S} h_k[n]$.
Since $S$ is time-invariant, we have $h_k[n] = h[n-k]$.
Together with the linearity, we have
\begin{align*}
    x[n] = \sum_{k=-\infty}^{\infty} x[k] \delta[n-k]
    \quad\xlongrightarrow{S}\quad y[n] &=
    \sum_{k=-\infty}^{\infty} x[k] h_k[n] \\
    &= \sum_{k=-\infty}^{\infty} x[k] h[n-k] \\
    &= x[n] \ast h[n]
\end{align*}

\begin{remark}
    In the impulse representation $x[n] \ast \delta[n]$,
    $x[n]$ becomes a constant sequence of scalers,
    and $\delta[n]$ becomes the one that actually carries the signal.
\end{remark}

The same for continuous-time:
\begin{align*}
    x(t) = \int_{-\infty}^{\infty} x(\tau) \delta(t-\tau) \dd{\tau}
    \quad\xlongrightarrow{S}\quad y(t) &=
    \int_{-\infty}^{\infty} x(\tau) h_\tau(t) \dd{\tau} \\
    &= \int_{-\infty}^{\infty} x(\tau) h(t-\tau) \dd{\tau} \\
    &= x(t) \ast h(t)
\end{align*}


\end{document}
